\documentclass[10pt,a4paper]{amsart}
\usepackage[margin=19mm]{geometry}
\usepackage{amsmath,amssymb,mathtools}
\usepackage[scr=rsfs,cal=euler]{mathalfa}
\usepackage{xfrac}
\usepackage[unicode,hidelinks,bookmarks=false]{hyperref}
\newcommand{\ie}{i.e.\ }
\newcommand{\cf}{cf.\ }
\newcommand{\resp}{respectively\ }
\newcommand{\Subsection}{\subsection}
\providecommand{\nobreakdash}{\nobreak\hbox{-}}
\setlength{\emergencystretch}{2em}
\multlinegap0pt
\usepackage{bm}
\usepackage{bbm}
\usepackage{dsfont}
\usepackage{xspace}
\usepackage{cancel}
\usepackage{mathtools}
\usepackage{amsthm}
\makeatletter
\@ifundefined{theorem}{\newtheorem{theorem}{Theorem}}{}
\@ifundefined{lemma}{\newtheorem{lemma}[theorem]{Lemma}}{}
\@ifundefined{remark}{\newtheorem{remark}[theorem]{Remark}}{}
\makeatother
\DeclareMathOperator*{\esssup}{ess\,sup}
\newcommand{\iprod}[1]{\langle #1 \rangle}
\newcommand{\pdiff}[2]{\frac{\partial #1}{\partial #2}}
\title[Intrinsic Finite Element Error Analysis on Manifolds with Re]{Intrinsic Finite Element Error Analysis on Manifolds with Regge Metrics, with Application to Calculating Connection Forms}
\author{Evan S. Gawlik, Jack McKee}
\date{Source: arXiv:2410.15579v3 (CC BY 4.0)}
\begin{document}
\maketitle

\noindent\emph{Source and licence.} Intrinsic Finite Element Error Analysis on Manifolds with Regge Metrics, with Application to Calculating Connection Forms. Version arXiv:2410.15579v3 (CC BY 4.0), distributed under the Creative Commons Attribution 4.0 International licence, as recorded in the arXiv licence metadata for this exact version and shown on the abstract page https://arxiv.org/abs/2410.15579v3. Full rights notice: licenses/arxiv-2410.15579-CC-BY-4.0.txt.\par\smallskip

\noindent\textit{Editorial scope.} Counted unit: the Theorem whose source label is \texttt{H1tracethrm} in the article (source0.tex lines 250--259, source label \texttt{H1tracethrm}), delivered together with the article's own complete original proof of it (source0.tex lines 265--331, one \texttt{proof} environment of 67 lines). No mathematics is written, repaired or re-derived: statement and proof are the article's own text line for line. The proof is detached from the statement in the source: the article states the result at lines 250--259 and prints its proof at lines 265--331, and the proof environment's own header names the statement, so the association is the article's own. Required context retained uncounted: L2\_equiv (lines 207--213); compare\_cov (lines 224--233); simplecomparecov (lines 237--242). Every definition, display and label of those lines is retained unchanged. Omitted from this package and not redistributed: 1--206, 214--223, 234--236, 243--249, 260--264, 332--920. Nothing inside a retained excerpt is omitted or altered: every retained body is delivered exactly as the article's lines are written, so no omission receipt of any kind accompanies this package. Citations: the retained excerpts carry no citation command at all, so no bibliography entry of the article is transcribed and no reference is redistributed. Reference adaptation: none was needed and none was made; every reference inside the delivered unit resolves inside it, so no unresolved reference or citation is produced. Printed slips kept literal: no printed word, symbol, hypothesis or number of the counted statement or of its proof was altered. Layout and class: the article's own class is \texttt{article} with packages including geometry, amsmath, amsfonts, amsthm, amssymb, xcolor, hyperref, biblatex, graphicx; the wrapper is the lane's pinned \texttt{amsart} editorial wrapper at 10pt on A4, which restates the theorem-like environments the retained text needs and adds the article's own macro definitions that the retained text needs (\texttt{\textbackslash esssup}, \texttt{\textbackslash iprod}, \texttt{\textbackslash pdiff}), with extra wrapper packages (bm, bbm, dsfont, xspace, cancel, mathtools). The delivered numbering is the unit's own. Assets: no retained span contains a figure environment, a graphics-inclusion command, a PDF-page inclusion, a table or a TikZ picture; no image file is copied into this package, the package declares no asset dependency and no image file is redistributed with it. Licence: version arXiv:2410.15579v3 of this article is distributed under the Creative Commons Attribution 4.0 International licence, as recorded in the arXiv licence metadata for this exact version and shown on the abstract page https://arxiv.org/abs/2410.15579v3. A full rights notice is delivered at licenses/arxiv-2410.15579-CC-BY-4.0.txt. Characters: every retained excerpt is pure ASCII, so the delivered engine, XeLaTeX with its default Latin Modern text font, reports no missing character.
\begin{theorem}\label{L2_equiv}
If $(\mathcal{T}_1,g_1)$ and $(\mathcal{T}_2,g_2)$ are two Regge metrics on the same manifold $M$, then
\begin{enumerate}
\item{$ \frac{1}{C_{g_1,g_2}(M)^n} \le D_{g_2,g_1}(M) \le C_{g_2,g_1}(M)^n$},
\item{$\displaystyle{\forall \alpha \in L^2\Omega^{p,k}(M,g_2), \;\|\alpha\|_{L^2(M,g_1)}^2 \le{n \choose k}n^pC_{g_2,g_1}(M)^{2p + 2k}D_{g_1,g_2}(M)\|\alpha\|_{L^2(M,g_2)}^2}$.}
\end{enumerate}
\end{theorem}

\begin{lemma}
\label{compare_cov}
Given two Regge metrics $g$ and $g'$ on a manifold $M$ and their Levi-Civita connections $\nabla$ and $\nabla'$, $\nabla - \nabla'$ is a bundle morphism $T^*(T \cap T')^p \otimes \Lambda^k(T \cap T') \to T^*(T \cap T')^{p+1} \otimes \Lambda^k(T \cap T')$ on each $T, T' \subseteq M$ where $g$ and $g'$ are smooth respectively. Furthermore, if $U \subseteq T \cap T'$ is measurable and contained in a coordinate chart with coordinates $\{x^i\}_{i = 1}^n$, then
\begin{align*}
\|\nabla - \nabla'\|_{\mathcal{L}(L^2\Omega^{p,k}(U,g),L^2\Omega^{p+1,k}(U,g))} \le \frac{3(k+p)n^\frac{7}{2}}{2}&\|G^{-1}\|_{L^\infty(U)}\|G\|_{L^\infty(U)}\\
&\cdot \big(\|G^{-1} - G'^{-1}\|_{L^\infty(U)}\|dG\|_{L^\infty(U,g)} \\
&\quad\quad + \|G'^{-1}\|_{L^\infty(U)}\|dG - dG'\|_{L^\infty(U,g)}\big),
\end{align*}
where $G_{ij} = \iprod{\pdiff{}{x^i},\pdiff{}{x^j}}_g$ is the coordinate expression for $g$ and $G'$ is defined similarly. For a matrix-valued one-form $A$, $\|A\|_{L^\infty(U,g)} := \esssup_{x \in U} \|A(x)\|_{2,g}$ and for a matrix-valued function, $\|A\|_{L^\infty(U)} := \esssup_{x \in U} \|A(x)\|_2$.
\end{lemma}

\begin{remark}
While the fact that the difference between Christoffel symbols is bounded by $\|G - G'\|$ and $\|dG - dG'\|$ is fairly obvious, this theorem makes explicit how they relate to bounds on the $L^2$ operator norm of $\nabla - \nabla'$. Additionally, if $g$ is piecewise flat, there exists a coordinate chart $U \subseteq T \cap T'$ around any point $x \in T \cap T'$ such that $G = \lambda I$ for a constant $\lambda > 0$ and $\|G'^{-1}\|_{L^\infty(U)} = 1$, yielding a much simpler, and intrinsic, expression:
\begin{equation}\label{simplecomparecov}\|\nabla - \nabla'\|_{\mathcal{L}(L^2\Omega^{p,k}(U,g),L^2\Omega^{p+1,k}(U,g))} \le \frac{3(k+p)n^\frac{7}{2}}{2}\|dG'\|_{L^\infty(U,g)} = \frac{3(k+p)n^\frac{7}{2}}{2}\|\nabla g'\|_{L^\infty(U,g)}\end{equation} 

Here $g'$ is understood to be a piecewise smooth $(0,2)$-tensor and $\nabla g'$ a $(0,3)$-tensor, and $\|\nabla g'\|_{L^\infty(U,g)} := \esssup_{x \in U}\|\nabla g'|_x\|_g$.
\end{remark}

\begin{theorem}[$H^1$ Trace Inequality for Riemannian Simplices]
\label{H1tracethrm}\;\\
Suppose $\alpha \in H^1\Omega^{p,k}(T,g)$, and assume that there exists an orientation-preserving diffeomorphism $f: \hat T \to T$ where $\hat T \subset \mathbb{R}^n$ is the standard simplex. Then $i_{\partial T}^*\alpha$ is well-defined and belongs to $L^2\Omega^{p,k}(\partial T,g)$, and there exists $\hat{C}$ depending only on $n$, $p$, and $k$ such that
\begin{align*}
\|&i_{\partial T}^*\alpha\|_{L^2(\partial T,g)} \\&\le \hat C  C_{g,\hat \delta}(T)^{k + p} C_{\hat \delta,g}(T)^{k+p + \frac{1}{2}}\sqrt{D_{g,\hat \delta}(T)D_{\hat \delta,g}(T)} \\
&\quad\cdot \bigg[\left(1 + C_{\hat \delta, g}(T)^{k+p + 1}C_{g,\hat \delta}(T)^{k+p + 1}\sqrt{D_{\hat \delta,g}(T)D_{g,\hat \delta}(T)} \|\hat \nabla g\|_{L^\infty(T,\hat \delta)}\right)\|\alpha\|_{L^2(T,g)} \\
&\quad\quad + C_{g,\hat \delta}(T)\|\nabla \alpha\|_{L^2(T,g)}\bigg],
\end{align*}
where $\hat \delta = f^{-*} \delta$ is the metric induced by $f$ and the standard Euclidean metric $\delta$ on $\hat T$, and $\hat\nabla$ is the associated connection.
\end{theorem}

\begin{proof}[Proof of Theorem \ref{H1tracethrm}]
We will take as given that there is a trace theorem for differential forms on $(\hat T, \delta)$. This is essentially because a form in this manifold is in $H^1$ if and only if each of its coefficients is in $H^1$ in the standard coordinates, and a trace theorem holds for each coefficient independently. Thus, setting $\hat \alpha = f^*\alpha$, $i_{\partial \hat T}^* \hat \alpha$ is well-defined and
\begin{equation}\label{basictraceeqn}
\|i_{\partial \hat T}^*\hat \alpha\|^2_{L^2(\partial \hat T,\delta)} \le \hat C \|\hat \alpha\|^2_{H^1(\hat T, \delta)}.
\end{equation}
Let $i_{\partial T}^*\alpha := f^{-*}i_{\partial \hat T}^*\hat \alpha$. This definition is consistent when $\alpha$ is continuous, because $f$ restricts to an embedding $e \hookrightarrow M$ for each facet $e \subset \partial T$. We will bound both sides of this inequality. Firstly, by Theorem \ref{L2_equiv},
\begin{align*}
%\|i_{\partial \hat T}^*\hat\alpha\|^2_{L^2(\partial \hat T,\delta)} &= \int_{\partial \hat T} \|f^*i_{\partial T}^* \alpha\|_\delta^2 dS_\delta = \int_{\partial T} \|i_{\partial T}^*\alpha\|_{f^{-*} \delta}^2 f^{-*}dS_\delta \\ &= \int_{\partial T} \|i_{\partial T}^*\alpha\|_{\hat \delta}^2 dS_{\hat \delta} = \|i_{\partial T}^*\alpha\|_{L^2(\partial T,\hat \delta)}^2 \\
\|i_{\partial \hat T}^*\hat\alpha\|^2_{L^2(\partial \hat T,\delta)} &= \|i_{\partial T}^*\alpha\|_{L^2(\partial T,\hat \delta)}^2 \\
&\ge \frac{1}{N_1^2C_{\hat \delta,g}(\partial T)^{2k + 2p}D_{g,\hat \delta}(\partial T)}\|i_{\partial T}^*\alpha\|_{L^2(\partial T,g)}^2,
\end{align*}
where $N_1$ is a constant depending only on $n$, $p$, and $k$.
Secondly,
\begin{align*}
%\|\hat{\alpha}\|^2_{H^1(\hat T, \delta)} &= \int_{\hat T} \|f^*\alpha\|^2_\delta dV_\delta + \int_{\hat T} \|\nabla_\delta f^*\alpha\|_\delta^2dV_\delta\\
%&= \int_{T} \|\alpha\|_{\hat \delta}^2 dV_{\hat \delta} + \int_T \|\hat\nabla\alpha\|^2_{\hat \delta}dV_{\hat \delta}\\
\|\hat{\alpha}\|^2_{H^1(\hat T, \delta)} &= \|\alpha\|^2_{L^2(T,\hat \delta)} + \|\hat\nabla\alpha\|^2_{L^2(T,\hat \delta)}\\
&\le N_2^2C_{g,\hat \delta}(T)^{2k + 2p}D_{\hat \delta,g}(T)\left(\|\alpha\|_{L^2(T,g)}^2 + C_{g,\hat \delta}(T)^2\|\hat\nabla\alpha\|_{L^2(T,g)}^2\right),
\end{align*}
where $N_2$ is another constant depending only on $n$, $p$, and $k$.
Plugging both of these inequalities into inequality (\ref{basictraceeqn}), we obtain
\begin{equation}\label{H1traceintermediate1}
\begin{split}
\|i_{\partial T}^*\alpha\|^2_{L^2(\partial T, g)} &\le \hat C N_1^2N_2^2C_{\hat \delta,g}(\partial T)^{2(k+p)}C_{g,\hat \delta}(T)^{2(k+p)}D_{\hat \delta,g}(T)D_{g,\hat \delta}(\partial T) \\ &\quad\cdot \left(\|\alpha\|_{L^2(T,g)}^2 + C_{g,\hat \delta}(T)^2\|\hat\nabla\alpha\|_{L^2(T,g)}^2\right).
\end{split}
\end{equation}
It is clearly the case that $C_{g_1,g_2}(\partial T) \le C_{g_1,g_2}(T)$ for any two continuous metrics $g_1, g_2$ on $T$. We can also bound $D_{g,\hat \delta}(\partial T)$ in terms of $D_{g,\hat \delta}(T)$. In a relatively open neighborhood $U \subset T$ of a point $x_0$ of $\partial T$ which is contained in a smooth component of $\partial T$, $dV_{g}$ can be written as $\theta^1\wedge\star_g\theta^1$, where $\|\theta^1\|_g = 1$ and $\ker \theta^1|_x = T_x\partial T$ for all $x \in U \cap \partial T$. Evidently, $i_{\partial T}^*\star_g \theta^1$ is the $g$ volume form on $U \cap \partial T$.

At each $x$ in $U$, we have the following inequality:
\begin{align*}
\sup_{W_1,\dots,W_n \in T_xT\backslash\{0\}}&\frac{|dV_{g}(W_1,\dots,W_n)|}{\|W_1\|_{\hat \delta}\dots\|W_n\|_{\hat \delta}} \\&\ge \sup_{W_2,\dots,W_n \in \ker \theta^1|_x\backslash\{0\}}\sup_{W_1 \notin \ker \theta^1|_x}\frac{|\theta^1(W_1)| \cdot |\star_g\theta^1(W_2,\dots,W_n)|}{\|W_1\|_{\hat \delta}\|W_2\|_{\hat \delta} \dots\|W_n\|_{\hat \delta}}\\
&= \|\theta^1|_x\|_{\hat \delta}\|\star_g\theta^1|_x\|_{\hat \delta}.
\end{align*}
Therefore, we have
\[
D_{g,\hat \delta}(U) \ge D_{g,\hat \delta}(U \cap \partial T)\inf_{x \in U}\|\theta^1|_x\|_{\hat \delta}  \ge \frac{D_{g,\hat \delta}(U \cap \partial T)}{C_{\hat \delta,g}(U)}. 
\]
This local inequality implies a global inequality
\[
D_{g,\hat \delta}(\partial T) \le C_{\hat \delta,g}(T)D_{g,\hat \delta}(T).
\]
Plugging this into inequality (\ref{H1traceintermediate1}), we derive
\begin{align*}
\|i_{\partial T}^*\alpha\|_{L^2(\partial T,g)}^2 \le \hat C^2 N_1^2N_2^2 &C_{\hat \delta,g}(T)^{2(k+p) + 1}C_{g,\hat \delta}(T)^{2(k+p)} \\&\cdot D_{\hat \delta,g}(T)D_{g,\hat \delta}(T)\left(\|\alpha\|_{L^2(T,g)}^2 + C_{g,\hat \delta}(T)^2\|\hat\nabla\alpha\|_{L^2(T,g)}^2\right),
\end{align*}
or, more conveniently,
\begin{align}
\|i_{\partial T}^*\alpha\|_{L^2(\partial T,g)} \le \hat C N_1N_2 & C_{\hat \delta,g}(T)^{k+p + \frac{1}{2}}C_{g,\hat \delta}(T)^{k+p} \notag \\
&\cdot \sqrt{D_{\hat \delta,g}(T)D_{g,\hat \delta}(T)}\left(\|\alpha\|_{L^2(T,g)} + C_{g,\hat \delta}(T)\|\hat\nabla\alpha\|_{L^2(T,g)}\right).\label{tracebounds}
\end{align}
Lastly, we can use Lemma \ref{compare_cov} to assert that 
\[
\|\hat\nabla \alpha\|_{L^2(T,g)} \le \|\nabla\alpha\|_{L^2(T,g)} + \|\hat\nabla - \nabla\|_{\mathcal{L}(L^2\Omega^{p,k}(T,g),L^2\Omega^{p+1,k}(T,g))}\|\alpha\|_{L^2(T,g)}
.\]

Since $\hat \delta$ is flat on $T$, we have by inequality \ref{simplecomparecov} that $\|\hat \nabla - \nabla\|_{\mathcal{L}(L^2\Omega^{p,k}(T,\hat \delta),L^2\Omega^{p+1,k}(T,\hat \delta)} \le \frac{3(k+p)n^\frac{7}{2}}{2}\|\hat \nabla g\|_{L^\infty(T,\hat \delta)}$. To apply this inequality, we'll use
\begin{align*}
&\|\hat \nabla - \nabla\|_{\mathcal{L}(L^2\Omega^{p,k}(T,g),L^2\Omega^{p+1,k}(T,g))} \\ 
&\le \frac{\sup_{\beta \in L^2\Omega^{p+1,k}(T,g) \backslash \{0\}} \frac{\|\beta\|_{L^2(T,g)}}{\|\beta\|_{L^2(T,\hat \delta)}}}{\inf_{\beta \in L^2\Omega^{p,k}(T,g) \backslash \{0\}}\frac{\|\beta\|_{L^2(T,g)}}{\|\beta\|_{L^2(T,\hat \delta)}}} \|\hat \nabla - \nabla\|_{\mathcal{L}(L^2\Omega^{p,k}(T,\hat \delta),L^2\Omega^{p+1,k}(T,\hat \delta))} \\
&\le \hat C_1\hat C_2 C_{\hat \delta,g}(T)^{k+p + 1}C_{g,\hat \delta}^{k+p}(T)\sqrt{D_{\hat \delta,g}(T)D_{g,\hat \delta}(T)}\|\hat \nabla - \nabla\|_{\mathcal{L}(L^2\Omega^{p,k}(T,\hat \delta),L^2\Omega^{p+1,k}(T,\hat \delta))} \\
&\le \hat C_1 \hat C_2 \frac{3(k+p)n^\frac{7}{2}}{2}C_{\hat \delta,g}(T)^{k+p+1}C_{g,\hat \delta}(T)^{k+p}\sqrt{D_{\hat \delta,g}(T)D_{g,\hat \delta}(T)} \|\hat \nabla g\|_{L^\infty(T,\hat \delta)},
\end{align*}
where $\hat C_1$ and $\hat C_2$ depend only on $n$, $p$, and $k$.

Plugging this into inequality \ref{tracebounds} and grouping together constant multiples, we obtain the desired bound.

\end{proof}

\end{document}
